\documentclass[11pt,a4paper]{article}

\usepackage[utf8]{inputenc}
\usepackage[T1]{fontenc}
\usepackage[french]{babel}
\usepackage{lmodern}
\usepackage{amsmath,amssymb}
\usepackage{bm}
\usepackage{booktabs,array}
\usepackage[margin=2.5cm]{geometry}
\usepackage{hyperref}
\emergencystretch=3em

\newcommand{\vm}{\bm{m}}
\newcommand{\vp}{\bm{p}}
\newcommand{\vs}{\bm{s}}
\newcommand{\vu}{\bm{u}}
\newcommand{\vg}{\bm{g}}
\newcommand{\vJ}{\bm{J}}
\newcommand{\vT}{\bm{T}}
\newcommand{\vH}{\bm{H}}
\newcommand{\ve}{\bm{e}}
\newcommand{\vb}{\bm{b}}
\newcommand{\code}[1]{\texttt{#1}}

\title{Couples adiabatique et non adiabatique\\
à partir de l'accumulation de spin dans FeeLLGood}
\author{}
\date{\today}

\begin{document}

\maketitle

\begin{abstract}
FeeLLGood ne calcule pas séparément les couples adiabatique et non adiabatique : le transfert de
spin entre dans l'équation de Landau--Lifshitz--Gilbert (LLG) par un seul champ effectif,
proportionnel à l'accumulation de spin $\vs$. Cette note montre comment extraire les deux
contributions des distributions $\vm$ et $\vs$ (et du potentiel $V$) écrites dans les fichiers
\code{.sol}. Elle donne aussi la prédiction analytique de la limite locale (Zhang--Li) avec les
paramètres du code, qui servira de référence au script de post-traitement, et la décomposition
\emph{field-like} / \emph{damping-like} utilisée quand un polariseur fixe impose la direction du
spin injecté.
\end{abstract}

\section{Le couple de transfert de spin dans FeeLLGood}

\subsection{Le problème de diffusion de spin}

À chaque pas de temps, l'aimantation étant figée, FeeLLGood résout l'état stationnaire de
l'accumulation de spin $\vs$ (A/m) :
\begin{equation}
  \partial_i Q_{i\alpha} + \frac{s_\alpha}{\tau_{sf}}
  + \frac{(\vs\times\vm)_\alpha}{\tau_{sd}} = 0,
  \label{eq:diffusion}
\end{equation}
avec le courant de spin (modèle de Zhang, Levy et Fert, $\beta=P$, $\beta'=P'$)
\begin{equation}
  Q_{i\alpha} = -P\,\frac{\mu_B}{e}\,J_i\,m_\alpha
  - D_0\left[\partial_i s_\alpha - PP'\,m_\alpha\,(\vm\cdot\partial_i\vs)\right],
  \qquad \vJ = -\sigma\nabla V .
  \label{eq:Q}
\end{equation}
Les indices latins désignent l'espace, les grecs la composante du spin, et les indices répétés
sont sommés. Les temps de relaxation sont
\begin{equation}
  \frac{1}{\tau_{sd}} = \frac{D_0}{l_{sd}^2},\qquad
  \frac{1}{\tau_{sf}} = \frac{D_0\,(1-PP')}{l_{sf}^2}
  \label{eq:taus}
\end{equation}
dans un ferromagnétique, et $1/\tau_{sf} = D_0/l_{sf}^2$ dans un métal normal
(\code{spinAcc::integrales}, \code{src/spinAccumulationSolver.cpp}). $P'$ vaut 0 par défaut
(diffusion isotrope) ; $P'=P$ est l'hypothèse $\beta=\beta'$, nécessaire à l'équivalence avec
Valet--Fert. La correspondance des notations avec le code et les
fichiers de réglages est donnée dans le tableau~\ref{tab:params}.

\begin{table}[ht]
\centering
\begin{tabular}{>{$}l<{$} l l l}
\toprule
\text{Symbole} & Signification & Unité & FeeLLGood \\
\midrule
\vm & aimantation réduite, $|\vm|=1$ & -- & \code{u} \\
\vs & accumulation de spin & A/m & \code{s} \\
V & potentiel électrique (avec \code{jn\_multiplier}) & V & \code{V} \\
\sigma & conductivité & S/m & \code{sigma} \\
P & taux de polarisation de la conductivité ($\beta$) & -- & \code{P} \\
P' & taux de polarisation de la diffusion ($\beta'$) & -- & \code{P\_prime} (défaut 0) \\
D_0 & constante de diffusion, $2\sigma/(e^2N_0)$ & m$^2$/s & calculée \\
l_{sd},\ l_{sf} & longueurs transverse et de spin-flip & m & \code{l\_sd}, \code{l\_sf} \\
M_s & aimantation à saturation & A/m & \code{Ms} \\
\gamma_0 & rapport gyromagnétique ($\mu_0|\gamma|$) & m/(A$\cdot$s) & \code{gamma0} \\
\bottomrule
\end{tabular}
\caption{Notations.}
\label{tab:params}
\end{table}

\subsection{Le champ effectif et le couple}

L'accumulation de spin entre dans LLG par le champ effectif (\code{spinAcc::prepareExtraField},
\code{Tetra::calc\_Hst})
\begin{equation}
  \vH_{st} = \frac{D_0}{l_{sd}^2\,\gamma_0 M_s}\,\vs
           = \frac{1}{\tau_{sd}\,\gamma_0 M_s}\,\vs ,
\end{equation}
ajouté au champ effectif micromagnétique. Le couple correspondant, en s$^{-1}$, dans
$\partial_t\vm = -\gamma_0\,\vm\times\vH_{\text{eff}} + \alpha\,\vm\times\partial_t\vm + \vT$, est
\begin{equation}
  \boxed{\;\vT = -\gamma_0\,\vm\times\vH_{st}
  = -\frac{1}{\tau_{sd}\,M_s}\,\vm\times\vs\;}
  \label{eq:T}
\end{equation}
Seule la partie transverse $\vs_\perp = \vs - (\vm\cdot\vs)\,\vm$ contribue ; $\vT$ est
perpendiculaire à $\vm$. Le couple~\eqref{eq:T} est \emph{le} couple de transfert de spin du
modèle : il contient les termes adiabatique et non adiabatique, et aussi les contributions non
locales (diffusion transverse) et celles des interfaces et des surfaces d'injection. Les sections
suivantes le décomposent.

\section{Décomposition dans le repère de la texture}

\subsection{Vitesse de dérive de spin}

On définit la vitesse de dérive de spin, dirigée dans le sens du mouvement des électrons :
\begin{equation}
  \vu = -\frac{P\,\mu_B}{e\,M_s}\,\vJ = \frac{P\,\mu_B\,\sigma}{e\,M_s}\,\nabla V ,
  \label{eq:u}
\end{equation}
et le vecteur
\begin{equation}
  \vg = (\vu\cdot\nabla)\,\vm ,\qquad g_\alpha = u_i\,\partial_i m_\alpha .
\end{equation}
Comme $|\vm|=1$, $\vm\cdot\vg = 0$ : $\vg$ est dans le plan tangent.

\subsection{Repère local et projections}

Là où $\vg\neq\bm{0}$, les vecteurs
\begin{equation}
  \ve_1 = \frac{\vg}{|\vg|},\qquad \ve_2 = \vm\times\ve_1
\end{equation}
forment une base orthonormée du plan perpendiculaire à $\vm$. On écrit
\begin{equation}
  \vT = \vT_{\text{ad}} + \vT_{\text{na}},\qquad
  \vT_{\text{ad}} = (\vT\cdot\ve_1)\,\ve_1,\qquad
  \vT_{\text{na}} = (\vT\cdot\ve_2)\,\ve_2 .
  \label{eq:decomp}
\end{equation}
Le couple adiabatique est porté par $(\vu\cdot\nabla)\vm$, le couple non adiabatique par
$\vm\times(\vu\cdot\nabla)\vm$. Avec $\ve_2\times\vm = \ve_1$ et $\ve_1\times\vm = -\ve_2$, les
projections de~\eqref{eq:T} s'expriment directement sur $\vs$ :
\begin{equation}
  \vT\cdot\ve_1 = \frac{\vs\cdot\ve_2}{\tau_{sd}\,M_s},\qquad
  \vT\cdot\ve_2 = -\frac{\vs\cdot\ve_1}{\tau_{sd}\,M_s}.
  \label{eq:proj}
\end{equation}
Le couple adiabatique vient donc de la composante de $\vs$ selon $\vm\times(\vu\cdot\nabla)\vm$,
et le couple non adiabatique de sa composante selon $(\vu\cdot\nabla)\vm$.

\subsection{Coefficients locaux}

Par comparaison avec la forme de Zhang--Li
\begin{equation}
  \vT_{\text{ZL}} = \eta\left[-(\vu\cdot\nabla)\vm + \beta\,\vm\times(\vu\cdot\nabla)\vm\right],
  \label{eq:ZL}
\end{equation}
on définit en chaque point où $|\vg|$ n'est pas négligeable une efficacité adiabatique et un
coefficient non adiabatique locaux :
\begin{equation}
  \boxed{\;
  \eta_{\text{loc}} = -\frac{\vT\cdot\ve_1}{|\vg|} = -\frac{\vs\cdot\ve_2}{\tau_{sd}M_s\,|\vg|},
  \qquad
  \beta_{\text{loc}} = -\frac{\vT\cdot\ve_2}{\vT\cdot\ve_1} = \frac{\vs\cdot\ve_1}{\vs\cdot\ve_2}
  \;}
  \label{eq:coefs}
\end{equation}
$\beta_{\text{loc}}$ ne dépend ni de $\tau_{sd}$ ni de $M_s$ : c'est le rapport des deux
composantes transverses de $\vs$ dans le repère de la texture.

\section{Prédiction analytique : la limite locale}
\label{sec:local}

\subsection{Dérivation}

On considère un ferromagnétique homogène traversé par un courant à divergence nulle, loin des
interfaces et des surfaces d'injection, et une texture qui varie lentement à l'échelle de
$l_{sd}$. On néglige alors les termes de diffusion ($D_0\,\partial_i\partial_i\vs$ et le terme en
$PP'$) devant les termes de relaxation. Avec $\partial_iJ_i = 0$, la divergence du terme polarisé
de~\eqref{eq:Q} donne $-P\,(\mu_B/e)\,(\vJ\cdot\nabla)\vm = M_s\,\vg$, et~\eqref{eq:diffusion}
devient
\begin{equation}
  \frac{\vs}{\tau_{sf}} + \frac{\vs\times\vm}{\tau_{sd}} = -M_s\,\vg .
\end{equation}
La composante selon $\vm$ donne $\vm\cdot\vs = 0$ à cet ordre ($\vm\cdot\vg = 0$). On pose
$\xi = \tau_{sd}/\tau_{sf}$ et on cherche $\vs = a\,\vg + b\,\vm\times\vg$. Avec
$\vg\times\vm = -\vm\times\vg$ et $(\vm\times\vg)\times\vm = \vg$ :
\begin{equation}
  \xi\,a + b = -\tau_{sd}M_s,\qquad \xi\,b - a = 0
  \quad\Longrightarrow\quad
  \vs = -\frac{\tau_{sd}M_s}{1+\xi^2}\left[\xi\,\vg + \vm\times\vg\right].
  \label{eq:s_local}
\end{equation}
Le couple~\eqref{eq:T}, avec $\vm\times(\vm\times\vg) = -\vg$, vaut
\begin{equation}
  \vT = \frac{1}{1+\xi^2}\left[-(\vu\cdot\nabla)\vm + \xi\,\vm\times(\vu\cdot\nabla)\vm\right],
\end{equation}
c'est-à-dire la forme~\eqref{eq:ZL} avec
\begin{equation}
  \boxed{\;
  \beta = \xi = \frac{\tau_{sd}}{\tau_{sf}} = (1-PP')\,\frac{l_{sd}^2}{l_{sf}^2},
  \qquad
  \eta = \frac{1}{1+\xi^2}
  \;}
  \label{eq:beta}
\end{equation}
Avec la valeur par défaut $P'=0$, $\beta = l_{sd}^2/l_{sf}^2$. Comme $\xi\ll 1$ en général, $\eta\simeq 1$ et le couple
adiabatique est le couple de Zhang--Li \og complet \fg{} $-(\vu\cdot\nabla)\vm$.

\subsection{Ordres de grandeur}

Pour $J = 10^{12}$~A/m$^2$, $P = 0{,}7$ et $M_s = 8\times10^5$~A/m :
$\mu_B/e = 5{,}79\times10^{-5}$~m$^2$/s et $|\vu| \simeq 51$~m/s. Avec $l_{sd} = 1$~nm et
$l_{sf} = 5$~nm, $\beta = 0{,}04$ pour $P'=0$ et $\beta = 0{,}51\times0{,}04 \simeq 0{,}020$ pour
$P'=P$.

\subsection{Validité}

La limite locale néglige la diffusion transverse $D_0\,\vs''$ devant la précession
$\vs\times\vm/\tau_{sd}$. Le premier ordre de la correction s'obtient en réinjectant
$l_{sd}^2\,\vs_0''$ (avec $\vs_0$ donné par~\eqref{eq:s_local}) dans le second membre. Pour une
paroi $\theta(z) = 2\arctan e^{(z-z_0)/\delta}$, $\vm = \sin\theta\,\hat{\bm x} +
\cos\theta\,\hat{\bm z}$, on a $\vg = u\,\theta'\,\bm e_\theta$ et
$\vm\times\vg = u\,\theta'\,\hat{\bm y}$, de direction fixe. Le terme dominant de
$\vs_0''$, $-K u\,\theta'''\,\hat{\bm y}$ avec $K = \tau_{sd}M_s/(1+\xi^2)$, produit par
précession la composante $l_{sd}^2 K u\,\theta'''$ selon $\bm e_\theta$ : le rapport des deux
composantes transverses de $\vs$ devient
\begin{equation}
  \beta_{\text{loc}}(z) = \xi - l_{sd}^2\,\frac{\theta'''}{\theta'}
  = \xi - \frac{l_{sd}^2}{\delta^2}\left(1 - 2\,\mathrm{sech}^2\frac{z-z_0}{\delta}\right),
  \qquad
  \bar\beta = \xi + \frac{l_{sd}^2}{3\delta^2},
  \label{eq:nonlocal}
\end{equation}
$\bar\beta$ étant la moyenne pondérée par $|\vg|^2\propto\mathrm{sech}^2$ de la
section~\ref{sec:num}. $\beta_{\text{loc}}$ n'est donc pas uniforme : plus grand au centre de la
paroi ($\xi + l_{sd}^2/\delta^2$), plus petit dans ses ailes. La correction est absolue, la même
pour tout $P'$ ; rapportée à $\xi$, elle vaut $(l_{sf}/\delta)^2/[3(1-PP')]$ : la limite locale
demande une paroi large devant $l_{sf}$, et pas seulement devant $l_{sd}$. Pour $l_{sf} = 5$~nm et
$\delta = 25$~nm, elle fait 1,3~\% de $\xi$ avec $P'=0$ et 2,6~\% avec $P'=P=0{,}7$.
\code{validation/wide\_wall} retrouve~\eqref{eq:nonlocal} à moins de 2~\% de la correction près,
pour $\delta$ = 12,5 et 25~nm. Quand $\delta$ s'approche de $l_{sd}$, les ordres suivants
interviennent et $\beta_{\text{loc}}$ peut largement dépasser $\xi$ (Claudio-Gonzalez, Thiaville et
Miltat, 2012 ; Petitjean, Luc et Waintal, 2012). C'est ce que la décomposition
numérique~\eqref{eq:decomp} permet de mesurer : l'écart entre $\beta_{\text{loc}}$
et~\eqref{eq:beta} quantifie les effets non locaux.

\section{Remarques}

\paragraph{Ce que la décomposition ne contient pas.} FeeLLGood résout l'état stationnaire de
$\vs$ : le terme $\partial_t\vs$, qui chez Zhang et Li renormalise le rapport gyromagnétique et
l'amortissement, est absent du modèle. La décomposition~\eqref{eq:decomp} porte donc sur le seul
couple~\eqref{eq:T}.

\paragraph{Interfaces et surfaces d'injection.} Près d'une surface où du spin est injecté
(conditions aux limites \code{uP}, \code{slonczewski} ou \code{tunnel}) ou d'une interface entre
régions, $\vs$ contient une partie qui ne vient pas de la texture : sa direction est imposée par
le spin injecté. La base $(\ve_1,\ve_2)$ reste définie si $\vg\neq\bm 0$, mais la séparation
adiabatique / non adiabatique n'y a plus le sens de Zhang--Li ; on utilise alors la décomposition
de la section~\ref{sec:polariseur}.

\paragraph{Partie longitudinale.} La composante $(\vm\cdot\vs)\,\vm$ ne crée aucun couple. Elle est
nulle dans la limite locale mais pas en général (interfaces, variations de $\sigma$, diffusion) :
il faut projeter, ce que font automatiquement~\eqref{eq:proj} et~\eqref{eq:coefs}.

\paragraph{Zones uniformes.} Là où $|\vg|$ est petit (domaines), $\ve_1$ est mal défini et
$\beta_{\text{loc}}$ est le rapport de deux quantités petites, dominé par l'erreur de
discrétisation. On masque ces points, par exemple $|\vg| < \varepsilon\,\max|\vg|$ avec
$\varepsilon$ de l'ordre de $10^{-2}$, ou on raisonne sur des quantités intégrées
(section~\ref{sec:num}).

\section{Polariseur fixe : couples field-like et damping-like}
\label{sec:polariseur}

Quand le spin est injecté par une couche de référence d'aimantation fixe $\vp$ (pile
\code{slonczewski} ou jonction \code{tunnel}), la base naturelle du plan perpendiculaire à $\vm$
est construite avec $\vp$ (elle existe si $\vm\nparallel\vp$) :
\begin{equation}
  \vb_1 = \frac{\vm\times\vp}{|\vm\times\vp|},\qquad
  \vb_2 = \frac{\vm\times(\vm\times\vp)}{|\vm\times\vp|} = \vm\times\vb_1 .
\end{equation}
Avec $(\vm\times\vs)\cdot(\vm\times\bm w) = \vs\cdot\bm w - (\vm\cdot\vs)(\vm\cdot\bm w)$ :
\begin{align}
  T_{\text{FL}} &= \vT\cdot\vb_1
    = -\frac{\vs\cdot\vp - (\vm\cdot\vs)(\vm\cdot\vp)}{\tau_{sd}\,M_s\,|\vm\times\vp|}, \\
  T_{\text{DL}} &= \vT\cdot\vb_2
    = -\frac{\vs\cdot(\vm\times\vp)}{\tau_{sd}\,M_s\,|\vm\times\vp|},
\end{align}
et $\vT = T_{\text{FL}}\,\vb_1 + T_{\text{DL}}\,\vb_2$. Le rapport $T_{\text{FL}}/T_{\text{DL}}$
est le rapport field-like / damping-like local ; il se compare à $G_i/G_r$ à l'interface. Comme
le couple décroît sur quelques $l_{sd}$ dans la couche libre, la quantité à comparer au modèle de
Slonczewski est le couple intégré sur la couche (section~\ref{sec:num}).

\section{Mise en œuvre numérique}
\label{sec:num}

\subsection{Données}

Avec \code{spin\_accumulation} activé, chaque ligne d'un fichier \code{.sol}
(\code{Mesh::mesh::savesol}, \code{src/save.cpp}) contient
\begin{center}
\code{idx \quad mx my mz \quad phi \quad sx sy sz \quad V}
\end{center}
aux nœuds du maillage, dans l'ordre des nœuds du \code{.msh}. Le $\vs$ écrit a été calculé avec
l'aimantation écrite sur la même ligne (\code{compute\_all} résout la diffusion de spin avant
\code{evolution()} puis \code{saver}). Le potentiel écrit inclut déjà \code{jn\_multiplier}
(\code{spinAcc::potential}). Les paramètres $P$, $\sigma$, $M_s$, $l_{sd}$, $l_{sf}$ et $N_0$ (pour
$D_0$) sont lus par région volumique dans le fichier de réglages.

\subsection{Calcul par tétraèdre}

Les champs sont P1 : dans un tétraèdre de nœuds $k=1\ldots4$ et de fonctions de forme $a_k$, les
gradients $\nabla a_k$ sont constants. Dans chaque tétraèdre magnétique :
\begin{enumerate}
  \item $\nabla V = \sum_k V_k\,\nabla a_k$, puis $\vu$ par~\eqref{eq:u} ;
  \item $\partial_i\vm = \sum_k \vm_k\,\partial_i a_k$ et $\vg = u_i\,\partial_i\vm$ ;
  \item $\vm$ et $\vs$ au barycentre (moyenne des nœuds), $\vm$ renormalisé ; on projette
    $\vg \leftarrow \vg - (\vm\cdot\vg)\,\vm$, l'interpolé P1 de $\vm$ n'étant pas unitaire ;
  \item $\vT$ par~\eqref{eq:T}, puis $\ve_1$, $\ve_2$, les projections~\eqref{eq:decomp} et les
    coefficients~\eqref{eq:coefs} (ou $\vb_1$, $\vb_2$ et $T_{\text{FL}}$, $T_{\text{DL}}$).
\end{enumerate}
On peut aussi évaluer $\vm$ et $\vs$ aux points de Gauss, comme le code le fait pour $\vH_{st}$ ;
au barycentre, l'écart est d'ordre $h^2$.

\subsection{Quantités intégrées}

Pour éviter les divisions par de petites quantités, on forme des moyennes pondérées par $|\vg|^2$
sur une région (une paroi, la couche libre) de volume $\Omega$ :
\begin{equation}
  \bar\eta = -\frac{\int_\Omega \vT\cdot\vg\,\mathrm dV}{\int_\Omega |\vg|^2\,\mathrm dV},
  \qquad
  \bar\beta\,\bar\eta = \frac{\int_\Omega \vT\cdot(\vm\times\vg)\,\mathrm dV}
  {\int_\Omega |\vg|^2\,\mathrm dV},
\end{equation}
qui sont les coefficients de la forme~\eqref{eq:ZL} ajustée au sens des moindres carrés sur
$\Omega$. Pour le polariseur, le moment cinétique transféré par unité de temps est
$\int_\Omega M_s\,\vT\,\mathrm dV$ (A$\cdot$m$^2$/s), dont on prend les composantes selon
$\vb_1$ et $\vb_2$ (avec $\vm$ uniforme dans une couche mince).

\subsection{Vérifications}

\begin{itemize}
  \item paroi large dans un fil uniforme : $\bar\beta$ et $\beta_{\text{loc}}$ doivent suivre la
    correction au premier ordre~\eqref{eq:nonlocal}, $\bar\eta$ tendre vers $1/(1+\xi^2)$ ;
    c'est le cas test \code{validation/wide\_wall} ;
  \item en faisant décroître $\delta$ vers $l_{sd}$, $\bar\beta$ montre le passage au régime non
    local ;
  \item convergence en maillage : la taille de maille doit résoudre $l_{sd}$ ;
  \item en faisant varier \code{P\_prime} de 0 à $P$, $\beta$ passe de $l_{sd}^2/l_{sf}^2$ à
    $(1-P^2)\,l_{sd}^2/l_{sf}^2$.
\end{itemize}

\section*{Références}

\begin{itemize}
  \item S. Zhang, P. M. Levy, A. Fert, Phys. Rev. Lett. \textbf{88}, 236601 (2002).
  \item S. Zhang, Z. Li, Phys. Rev. Lett. \textbf{93}, 127204 (2004).
  \item A. Thiaville, Y. Nakatani, J. Miltat, Y. Suzuki, Europhys. Lett. \textbf{69}, 990 (2005).
  \item D. Claudio-Gonzalez, A. Thiaville, J. Miltat, Phys. Rev. Lett. \textbf{108}, 227208 (2012).
  \item C. Petitjean, D. Luc, X. Waintal, Phys. Rev. Lett. \textbf{109}, 117204 (2012).
\end{itemize}

\end{document}
